\section{Strict interlacing through the first four-cycle transition}
\label{sec:three-cycle}
Fix $k\ge3$ and write $P_n=E_{n,k}/t$. The two-cycle theorem of
Section~\ref{sec:two-cycle-insertion} treats $k\le n\le 3k-2$.
We extend the range to $k\le n\le4k-1$. In particular the last new
row, at size $4k$, includes configurations with four cycles.

\paragraph{Theorem.}
For $3k-1\le n<4k$, at every zero $x$ of $P_n$ one has
\[
\left|\binom n{k-1}A_{k-1}(x) E_{n+1-k,k}(x)\right|
  < |x(1-x)P_n'(x)|.
\]
The $n-2$ zeros of $P_n$ and the $n-1$ zeros of $P_{n+1}$ are
simple and negative, and strictly alternate. The new row contributes
both outermost zeros.

\paragraph{Large sizes.}
The generating function in the separated-pole argument is a product
over odd poles. Write $t=-\exp(-\pi\cot\theta)$,
$S=\sin\theta$, $z_p=(\cos\theta+ip S)^{-1}$ for $p$ odd, and
$F_j=(\pi/S)^j E_{j,k}(t)/(j!(1-t)^j)$. A pole contributes either
its constant term or at least $k$ units of degree. In the old and
residual rows the nonconstant terms can consequently occupy at most
three poles, even at $n=4k-1$. Partial fractions give the single-pole
main term $d_n\sum_p z_p^n$, where
$d_j=[u^j]\exp(\sum_{s\ge k}u^s/s)$.
Integer phase samples $\theta=j\pi/n$ have sign $(-1)^j$
on the middle interval. The bounds on the differentiated remainder
and on the insertion term are given with a rational checker in the
\href{https://github.com/ProgramComputer/open-math-research/blob/main/topics/long-cycle-excedances/three-cycle-2200.md}{three-cycle estimates}.
Their middle comparison holds for $n\ge192$, $kS\ge14$ and
$nS^2\ge149/100$. Writing $B_k=\sum_p z_p^k$ and $m=n+1-k$,
at every middle zero the two comparable magnitudes satisfy
\[
 k|F_n'|/n>19/100, \qquad k|B_k F_m|/(n S)<16/100 .
\]
The common normalization gives the claimed strict inequality. For
$n\ge2200$ the first $\lfloor .39\sqrt{n-1}\rfloor$ ranks on each
side are covered by a two-term Rouch\'e comparison. Its derivative
contribution exceeds $.63$ and its lag is below $.58$ in the same
scaling. The first middle sample beyond these ranks lies inside
the middle interval. Thus the sample signs and reciprocal reflection
exhaust the old degree.

\paragraph{Intermediate sizes.}
For $192\le n\le2199$ take $\theta_*=20\pi/n$. On this range
$k>n/4$ and $\sin u\ge u(1-u^2/6)$ up to $20\pi/192$ give
$k\sin\theta_*>14$ and $n\sin^2\theta_*>149/100$. There are
$n-40$ sample gaps between integer phases $20$ and $n-20$.
It remains to certify nineteen old roots on either side.

Let $I=t(1-t)E'_{n,k}+ntE_{n,k}$ and
$J=\binom n{k-1}t A_{k-1}E_{m,k}$. Their coefficients are
nonnegative, so the absolute coefficients of $I-J$ are bounded by
those of $I+J=E_{n+1,k}$. The decreasing-run encoding bounds
$[t^p]E_{N,k}$ by $p^N$. For $t=-x$, $x>0$, discarding powers
of degree $48$ and higher from $P_n$ or either next pencil
$(I\pm J)/t$ therefore costs at most
\[
 \frac{49^{n+1} x^{48}}{1-x(50/49)^{n+1}}
\]
when the denominator is positive. The verifier bounds both the
prefix evaluation and this remainder with outward-rounded real
balls. It accepts only strict signs. Between zero and phase~20
it isolates nineteen sign changes and refines the brackets until
both next pencils have the required signs at both ends. Each
middle zero has the required sign by the analytic inequality.
The intervening gaps, reflection, and the two outer intervals
account for all $n-1$ new zeros. Hence no pencil zero lies inside
an old bracket. This argument also proves the strict lag inequality
at the exterior old roots.

\paragraph{Small sizes and reproduction.}
For the $1{,}595$ pairs below size $192$, whole-polynomial
reconstructions and dyadic endpoint signs give the two strict signs
at $198{,}819$ old roots. In the intermediate range the calculation
uses $200{,}800$ pairs and $3{,}815{,}200$ exterior brackets before
reflection. The coefficient reconstruction, interval checks and
coverage records are available with the
\href{https://github.com/ProgramComputer/open-math-research/blob/main/topics/long-cycle-excedances/three-cycle-complete.md}{proof and reproduction instructions}.
The statement at all sizes follows by combining the three ranges.
Together with the known uniform upper cutoff and the all-size result
at $k=3$, a remaining counterexample must have $k\ge4$ and
$4k\le n<10^{20}$.
